
\documentclass[11pt,a4paper]{article}

\usepackage[T1]{fontenc}
\usepackage[utf8]{inputenc}
\usepackage{lmodern}
\usepackage[a4paper,margin=2.65cm]{geometry}
\usepackage{amsmath,amssymb,amsthm,mathtools}
\usepackage{enumitem}
\usepackage{hyperref}
\usepackage{xcolor}
\usepackage{array}
\usepackage{booktabs}
\usepackage{microtype}
\usepackage{longtable}

\hypersetup{
  colorlinks=true,
  linkcolor=blue!55!black,
  citecolor=blue!55!black,
  urlcolor=blue!55!black
}

\setlength{\parindent}{0pt}
\setlength{\parskip}{0.55em}
\setlist[itemize]{leftmargin=2.2em,itemsep=0.25em,topsep=0.35em}
\setlist[enumerate]{leftmargin=2.4em,itemsep=0.25em,topsep=0.35em}
\renewcommand{\arraystretch}{1.18}

\newcommand{\A}{\mathbb A}
\newcommand{\Z}{\mathbb Z}
\newcommand{\Q}{\mathbb Q}
\newcommand{\F}{\mathbb F}
\newcommand{\R}{\mathrm{R}}
\newcommand{\Gm}{\mathbb G_m}
\newcommand{\Pone}{\mathbb P^1}
\newcommand{\Aone}{\mathbb A^1}
\newcommand{\Spec}{\mathrm{Spec}}
\newcommand{\Sp}{\mathrm{Sp}}
\newcommand{\SH}{\mathrm{SH}}
\newcommand{\Sm}{\mathrm{Sm}}
\newcommand{\Sch}{\mathrm{Sch}}
\newcommand{\CAlg}{\mathrm{CAlg}}
\newcommand{\Gr}{\mathrm{Gr}}
\newcommand{\Sh}{\mathrm{Sh}}
\newcommand{\D}{\mathrm{D}}
\newcommand{\Nis}{\mathrm{Nis}}
\newcommand{\et}{\mathrm{\acute et}}
\newcommand{\cdh}{\mathrm{cdh}}
\newcommand{\BL}{\mathcal{BL}}
\newcommand{\BLp}{\mathcal{BL}^{p}}
\newcommand{\syn}{\mathrm{syn}}
\newcommand{\pbf}{\mathrm{pbf}}
\newcommand{\KH}{\mathrm{KH}}
\newcommand{\KGL}{\mathrm{KGL}}
\newcommand{\kgl}{\mathrm{kgl}}
\newcommand{\HZ}{\mathrm{HZ}}
\newcommand{\HF}{\mathrm{HF}}
\newcommand{\TC}{\mathrm{TC}}
\newcommand{\Pic}{\mathrm{Pic}}
\newcommand{\Fil}{\mathrm{Fil}}
\newcommand{\fib}{\mathrm{fib}}
\newcommand{\cofib}{\mathrm{cofib}}
\newcommand{\colim}{\operatorname*{colim}}
\newcommand{\limi}{\operatorname*{lim}}
\newcommand{\map}{\mathrm{map}}
\newcommand{\Map}{\mathrm{Map}}
\newcommand{\id}{\mathrm{id}}
\newcommand{\1}{\mathbf 1}
\newcommand{\ot}{\otimes}
\newcommand{\sm}{\mathrm{sm}}
\newcommand{\h}{\mathrm{h}}
\newcommand{\sh}{\mathrm{sh}}
\newcommand{\op}{\mathrm{op}}
\newcommand{\Mod}{\mathrm{Mod}}
\newcommand{\cA}{c^{A}_{1}}
\newcommand{\cBL}{c^{\BL}_{1}}
\newcommand{\leqj}{\tau^{\leq j}}
\newcommand{\gtj}{\tau^{>j}}
\newcommand{\leqstar}{\tau^{\leq \star}}
\newcommand{\slic}{\mathrm{slice}}
\newcommand{\veff}{\mathrm{veff}}
\newcommand{\eff}{\mathrm{eff}}
\newcommand{\qcqs}{\mathrm{qcqs}}
\newcommand{\Lnis}{L_{\Nis}}
\newcommand{\Let}{L_{\et}}
\newcommand{\Laone}{L_{\Aone}}
\newcommand{\Lcdh}{L_{\cdh}}
\newcommand{\Lan}{L^{\sm}}
\newcommand{\LL}{L}
\newcommand{\ZBL}{\Z_p(\star)^{\BL}}
\newcommand{\Zsyn}{\Z_p(\star)_{\syn}}
\newcommand{\ZA}{\Z(\star)^A}
\newcommand{\QCoh}{\mathrm{QCoh}}
\newcommand{\B}{\mathrm{B}}
\newcommand{\Tot}{\mathrm{Tot}}
\newcommand{\calO}{\mathcal O}
\newcommand{\calF}{\mathcal F}
\newcommand{\calC}{\mathcal C}
\newcommand{\calD}{\mathcal D}
\newcommand{\calE}{\mathcal E}
\newcommand{\calH}{\mathcal H}
\newcommand{\calP}{\mathcal P}
\newcommand{\calT}{\mathcal T}
\newcommand{\calU}{\mathcal U}
\newcommand{\calV}{\mathcal V}
\newcommand{\bKM}{\widehat K^M}

\newtheoremstyle{bemstyle}
  {0.7em}{0.7em}{\itshape}{} {\bfseries}{.}{0.5em}{}
\theoremstyle{bemstyle}
\newtheorem{theorem}{Theorem}[section]
\newtheorem{proposition}[theorem]{Proposition}
\newtheorem{lemma}[theorem]{Lemma}
\newtheorem{corollary}[theorem]{Corollary}
\theoremstyle{definition}
\newtheorem{definition}[theorem]{Definition}
\newtheorem{remark}[theorem]{Remark}

\title{Beilinson--Lichtenbaum over Dedekind Domains}
\author{Andrea Taccani}
\date{23 June 2026}

\begin{document}
\maketitle

\begin{abstract}
These notes explain the argument of Bachmann--Elmanto--Morrow, \S\S6.2--6.3, which passes from the Beilinson--Lichtenbaum truncation of syntomic cohomology to a motivic spectrum, proves its base-change property over fields and mixed-characteristic Dedekind domains, and deduces the zero-slice comparison and the Beilinson--Lichtenbaum equivalence for smooth schemes over such bases. The theorem and proposition numbers in the text refer to that paper.
\end{abstract}

\tableofcontents

\section{Goal of the talk}

The purpose of the talk is to explain how one passes from syntomic cohomology to a motivic spectrum and then transports field-level results to mixed-characteristic Dedekind domains. The logical chain is:
\[
\begin{gathered}
\text{syntomic cohomology}
\Rightarrow \Z_p(\star)^{\BL}
\Rightarrow \HZ^{\BL}_{p,B}
\Rightarrow \text{base change}\\[0.35em]
\Rightarrow s^0(\1_B)^\wedge_p\simeq \HZ^{\BL}_{p,B}
\Rightarrow \text{BL comparison on }\Sm_B.
\end{gathered}
\]

\section{BEM \S6.2: The Beilinson--Lichtenbaum motivic spectrum and base change}

\subsection{Prerequisites to recall at the beginning of the talk}

\subsubsection{Beilinson--Lichtenbaum cohomology}

From \S5, for each prime \(p\) and integer \(j\), BEM have syntomic cohomology
\[
  \Z_p(j)^{\syn} : \Sch^{\mathrm{qcqs},op}\longrightarrow \D(\Z).
\]
The Beilinson--Lichtenbaum version is obtained by truncating syntomic cohomology in degrees \(\leq j\), and then Nisnevich sheafifying. Modulo \(p\), the definition is
\[
  \F_p(j)^{\BL}:=L_{\Nis}\,\tau^{\leq j}\F_p(j)^{\syn}.
\]
The \(p\)-adic version is
\[
  \Z_p(j)^{\BL}:=\varprojlim_{r} L_{\Nis}\,\tau^{\leq j}(\Z_p(j)^{\syn}/p^r).
\]
The truncation is the Beilinson--Lichtenbaum truncation: in weight \(j\), one keeps cohomological degrees at most \(j\). 

For fields and mixed-characteristic Dedekind domains, Theorem 5.47 says that \(\Z_p(\star)^{\BL}\), restricted to smooth schemes over the base, is \(\Aone\)-invariant and satisfies the \(\Pone\)-bundle formula. 

\subsubsection{The motivic Eilenberg--MacLane construction}

BEM have a functor, denoted here by
\[
  H:\CAlg\bigl(\Gr\Sh_{\Nis,\Aone}(\Sm_B,\Sp)_{c,\pbf}\bigr)
  \longrightarrow \CAlg(\SH(B)).
\]
The input is a graded \(\Aone\)-invariant Nisnevich sheaf equipped with multiplication, a first Chern class, and the \(\Pone\)-bundle formula. The output is a motivic ring spectrum in \(\SH(B)\).

The defining compatibility is:
\[
  (HF)(j)[2j]\simeq F^j.
\]
Equivalently, if \(E=HF\), then on smooth \(B\)-schemes \(X\) one has
\[
  \operatorname{map}_{\SH(B)}(M_B(X),E\otimes T_B^{\otimes j}[-2j])\simeq F^j(X).
\]
This is the bridge from a cohomology theory to an object of motivic homotopy theory.

\subsubsection{Why base change is non-formal}

One might naively expect that if a cohomology theory is natural in the base, then the associated motivic spectrum automatically commutes with base change. This is false as stated. The construction \(H\) is compatible with base change only after passing through the correct category of \(\Aone\)-invariant sheaves satisfying the \(\Pone\)-bundle formula. Therefore BEM must prove that pulling \(\Z_p(\star)^{\BL}\) from \(B\) to \(k\) gives the same object as \(\Z_p(\star)^{\BL}\) constructed directly over \(k\), at least modulo \(p\). That is what Proposition 6.16 supplies.

\subsection{Construction of the spectrum \texorpdfstring{\(\HZ^{\BL}_{p,B}\)}{HZ BL p,B}}

Let \(B\) be either a field or a mixed-characteristic Dedekind domain. By Theorem 5.47, the graded object
\[
  \Z_p(\star)^{\BL}[2\star]\big|_{\Sm_B}
\]
is an \(E_\infty\)-algebra in graded \(\Aone\)-invariant Nisnevich sheaves on \(\Sm_B\), equipped with the first Chern class
\[
  \cBL(\mathcal O(1))\in H^2_{\BL}(\Pone_B,\Z_p(1)),
\]
and satisfying the \(\Pone\)-bundle formula. Therefore it defines an object
\[
  \Z^{\BL}_{p,B}\in
  \CAlg\bigl(\Gr\Sh_{\Nis,\Aone}(\Sm_B,\Sp)_{c,\pbf}\bigr).
\]
Applying the motivic Eilenberg--MacLane functor gives the motivic spectrum
\[
  \HZ^{\BL}_{p,B}:=H(\Z^{\BL}_{p,B})\in\CAlg(\SH(B)).
\]

By construction, \(\HZ^{\BL}_{p,B}\) represents Beilinson--Lichtenbaum cohomology on smooth \(B\)-schemes. Namely, for each \(j\in\Z\),
\[
  \HZ^{\BL}_{p,B}(j)(X)
  :=\operatorname{map}_{\SH(B)}(M_B(X),\HZ^{\BL}_{p,B}\otimes T_B^{\otimes j}[-2j])
  \simeq \Z_p(j)^{\BL}(X)
\]
for \(X\in\Sm_B\). 

There is also a mod-\(p\) version. Applying \(H\) to
\[
  \F_p(\star)^{\BL}[2\star]
\]
after reducing the first Chern class modulo \(p\), one gets
\[
  \HF^{\BL}_{p,B}\in\CAlg(\SH(B)).
\]
The paper notes that this comes equipped with an equivalence of \(\1_B\)-pointed motivic spectra
\[
  \HF^{\BL}_{p,B}\simeq \HZ^{\BL}_{p,B}/p.
\]
In particular, \(\HZ^{\BL}_{p,B}/p\) inherits an \(E_\infty\)-algebra structure.

\begin{remark}[What this construction is doing]
This construction converts a concrete cohomology theory into a motivic spectrum. The concrete object is \(\Z_p(j)^{\BL}\), built by truncating syntomic cohomology. The motivic object is \(\HZ^{\BL}_{p,B}\). The reason this matters is that slices, base change, and comparison with \(s^0(\1_B)\) are statements internal to \(\SH(B)\), not merely statements about complexes of abelian groups.
\end{remark}

\subsection{Proposition 6.16: the technical input}

\begin{proposition}[BEM Proposition 6.16]
Let \(j\in\Z\).
\begin{enumerate}
  \item Rigidity:  For any ring \(R\) and henselian ideal \(I\subset R\), the relative syntomic cohomology
  \[
    \F_p(j)^{\syn}(R,I):=\fib\bigl(\F_p(j)^{\syn}(R)\to \F_p(j)^{\syn}(R/I)\bigr)
  \]
  is supported in cohomological degrees \(\leq j\).
  \item Left Kan extended: For any base ring \(B\), the functor
  \[
    \tau^{\leq j}\F_p(j)^{\syn}:\CAlg_B\longrightarrow \D(\Z)
  \]
  is left Kan extended from smooth \(B\)-algebras.
\end{enumerate}
\end{proposition}

\subsubsection{Meaning of the statement}

The first assertion is not saying that all of syntomic cohomology is invariant under henselian quotients. It says something slightly weaker and suited to the Beilinson--Lichtenbaum truncation: the relative term has no cohomology above degree \(j\). Consequently, after applying \(\tau^{>j}\), the map
\[
  \F_p(j)^{\syn}(R)\longrightarrow \F_p(j)^{\syn}(R/I)
\]
becomes an equivalence. Thus the ``error'' of passing from \(R\) to \(R/I\) lives only in the degrees which the BL truncation keeps.

The second assertion is the input needed for base change. It says that \(\tau^{\leq j}\F_p(j)^{\syn}\) is determined by its restriction to smooth algebras. Combined with BEM Proposition 2.23 and Nisnevich sheafification, this identifies the mod-\(p\) pullback of the Beilinson--Lichtenbaum theory from \(B\) to \(k\) with the theory constructed directly over \(k\).

\subsubsection{Proof of Proposition 6.16(1)}

Let \(I\subset R\) be henselian. We must show
\[
  H^n\F_p(j)^{\syn}(R,I)=0\qquad (n>j).
\]
The proof proceeds by reductions until one can use the known rigidity theorem for syntomic cohomology of \(p\)-complete rings.

First, by Corollary 5.22, one may replace \(R\) by a \(p\)-henselian object for the purpose of the relative calculation. This reduction isolates the part of the ring relevant to syntomic cohomology at \(p\).

Second, using Proposition 5.20, syntomic cohomology modulo \(p\) is left Kan extended from smooth algebras and commutes with the required filtered colimits. Thus one may reduce to the Noetherian case. This is important because completions behave cleanly for Noetherian rings.

Third, Lemma 5.18(2), applied to \(R\) and to \(R/I\), identifies the relative syntomic term with the fibre
\[
  \fib\Bigl(
    \F_p(j)^{\syn}(R^\wedge_p)
    \longrightarrow
    \F_p(j)^{\syn}((R/I)^\wedge_p)
  \Bigr).
\]
Because \(R\) is now Noetherian, completion commutes with quotienting by \(I\):
\[
  (R/I)^\wedge_p\simeq R^\wedge_p/IR^\wedge_p.
\]
Moreover, \(IR^\wedge_p\) is a henselian ideal of \(R^\wedge_p\). The reason is that henselianity descends to the mod-\(p\) quotient
\[
  R^\wedge_p/pR^\wedge_p\simeq R/pR,
\]
where the \(I\) generates a henselian ideal.

Now one applies the rigidity theorem for syntomic cohomology of \(p\)-complete rings. It says that for a \(p\)-complete ring and a henselian ideal, the relative syntomic cohomology in weight \(j\) is concentrated in degrees \(\leq j\). This proves part (1).

\subsubsection{Proof of Proposition 6.16(2)}

Part (1) implies that the functor
\[
  \tau^{>j}\F_p(j)^{\syn}:\CAlg_B\longrightarrow \D(\Z)
\]
is rigid in the sense of BEM Definition 2.24: for every henselian ideal \(I\subset R\), the map
\[
  \tau^{>j}\F_p(j)^{\syn}(R)
  \longrightarrow
  \tau^{>j}\F_p(j)^{\syn}(R/I)
\]
is an equivalence. Indeed, the fibre of this map is \(\tau^{>j}\F_p(j)^{\syn}(R,I)\), and part (1) says this fibre vanishes.

By BEM Lemma 2.25, a finitary rigid functor is left Kan extended from smooth algebras. Hence \(\tau^{>j}\F_p(j)^{\syn}\) is left Kan extended from smooth \(B\)-algebras.

On the other hand, Proposition 5.20(2) says that \(\F_p(j)^{\syn}\) itself is left Kan extended from smooth \(B\)-algebras. Since there is a fibre sequence\footnote{This is the canonical truncation triangle for the standard \(t\)-structure on \(\D(\Z)\): for every \(K\in\D(\Z)\) there is a functorial fibre--cofibre sequence \(\tau^{\leq j}K\to K\to\tau^{\geq j+1}K\). Here \(\tau^{>j}=\tau^{\geq j+1}\), and the sequence is applied objectwise. See Lurie, \emph{Higher Algebra}, \S1.2.1 (t-structures).}
\[
  \tau^{\leq j}\F_p(j)^{\syn}
  \longrightarrow
  \F_p(j)^{\syn}
  \longrightarrow
  \tau^{>j}\F_p(j)^{\syn},
\]
and the two right-hand terms are left Kan extended from smooth algebras, the left-hand term is also left Kan extended from smooth algebras. This proves part (2).


\subsection{Theorem 6.17: base change for \texorpdfstring{\(\HZ^{\BL}_{p,B}\)}{HZ BL}}

\begin{theorem}[BEM Theorem 6.17]
Let \(B\) be a mixed-characteristic Dedekind domain or a field. Let \(k\) also be a mixed-characteristic Dedekind domain or a field, and let
\[
  i:B\longrightarrow k
\]
be a map of rings. Then there is an equivalence of \(E_\infty\)-algebras
\[
  (i^*\HZ^{\BL}_{p,B})^\wedge_p\simeq \HZ^{\BL}_{p,k}
\]
in \(\SH(k)\).
\end{theorem}

The important case is when \(B\) is a mixed-characteristic Dedekind domain and \(k\) is one of its residue fields or its fraction field. This is the situation needed in \S6.3, where one wants to reduce a statement over \(B\) to field-valued points.

\subsubsection{The compatibility diagram}

The proof uses the compatibility of the motivic Eilenberg--MacLane functor \(H\) with pullback. Schematically, the relevant diagram is
\[
\begin{array}{ccc}
\CAlg(\Gr\Sh_{\Nis,\Aone}(\Sm_k,\Sp)_{c,\pbf}) & \xrightarrow{\ H\ } & \CAlg(\SH(k)) \\
\uparrow L_{\pbf} && \uparrow i^* \\
\CAlg(\Gr\Sh_{\Nis,\Aone}(\Sm_k,\Sp)_{c}) && \\
\uparrow L_{\Nis,\Aone}L^{\sm} && \\
\CAlg(\Gr\Sh_{\Nis,\Aone}(\Sm_B,\Sp)_{c,\pbf}) & \xrightarrow{\ H\ } & \CAlg(\SH(B)).
\end{array}
\]
Read this diagram from bottom left to top right. Starting with the BL cohomology theory over \(B\), first pull it to smooth \(k\)-schemes by left Kan extension and sheafification; then impose \(\Aone\)-invariance and the projective bundle formula; then apply \(H\). This agrees with first applying \(H\) over \(B\) and then pulling the resulting motivic spectrum to \(k\).

\subsubsection{Proof of Theorem 6.17}

Start with the object
\[
  \Z^{\BL}_{p,B}\in\CAlg(\Gr\Sh_{\Nis,\Aone}(\Sm_B,\Sp)_{c,\pbf}).
\]
By naturality of Beilinson--Lichtenbaum cohomology there is a comparison map on the sheaf side
\[
  L_{\Nis}L^{\sm}\Z^{\BL}_{p,B}
  \longrightarrow
  \Z^{\BL}_{p,k}.
\]
Here \(L^{\sm}\) means left Kan extension from smooth \(B\)-schemes to smooth \(k\)-schemes.

Modulo \(p\), this comparison map is an equivalence. The reason is exactly Proposition 6.16(2), combined with BEM Proposition 2.23. Indeed, modulo \(p\), the BL theory is
\[
  \F_p(j)^{\BL}=L_{\Nis}\tau^{\leq j}\F_p(j)^{\syn},
\]
and Proposition 6.16 says that \(\tau^{\leq j}\F_p(j)^{\syn}\) is left Kan extended from smooth algebras over any base. Proposition 2.23 then says that these left Kan extensions commute with the base-change operation under consideration. Hence the pulled-back BL sheaf over \(B\) agrees modulo \(p\) with the BL sheaf constructed directly over \(k\).

The target \(\Z^{\BL}_{p,k}\) is \(\Aone\)-invariant and satisfies the \(\Pone\)-bundle formula by Theorem 5.47. Therefore the comparison map factors through the operation which imposes these properties:
\[
  L_{\pbf}L_{\Nis,\Aone}L^{\sm}\Z^{\BL}_{p,B}
  \longrightarrow
  \Z^{\BL}_{p,k}.
\]
This map remains an equivalence modulo \(p\). Informally: imposing \(\Aone\)-invariance and the \(\Pone\)-bundle formula does not alter the comparison modulo \(p\), because the target already has those properties and the domain was already correct modulo \(p\).

Now apply the motivic Eilenberg--MacLane functor \(H\). By the compatibility diagram above, this produces a map of motivic \(E_\infty\)-algebras
\[
  i^*\HZ^{\BL}_{p,B}
  \longrightarrow
  \HZ^{\BL}_{p,k}
\]
in \(\SH(k)\). The map is an equivalence modulo \(p\). Since \(\HZ^{\BL}_{p,k}\) is \(p\)-complete, passing to \(p\)-completion gives
\[
  (i^*\HZ^{\BL}_{p,B})^\wedge_p\simeq \HZ^{\BL}_{p,k}.
\]
This proves Theorem 6.17.

\begin{remark}[Canonicity]
At this point, we are not yet claiming a canonical equivalence. After Proposition 6.18, the object \(\HZ^{\BL}_{p,B}\) is identified with \(s^0(\1_B)^\wedge_p\), and that identification forces uniqueness. 
\end{remark}

The subsequent section uses base change to compare \(\HZ^{\BL}_{p,B}\) with the zero slice of the motivic sphere and with the zero slice of \(\KGL\). The chain of dependence is:
\[
\begin{array}{c}
\text{Theorem 5.47: BL cohomology has }\Aone\text{-invariance and PBF over fields/Dedekind domains}\\[0.4em]
\Downarrow\\[0.4em]
\text{Definition of }\HZ^{\BL}_{p,B}\text{ via }H\\[0.4em]
\Downarrow\\[0.4em]
\begin{gathered}
\text{Proposition 6.16: truncated syntomic cohomology}\\[-0.1em]
\text{is left Kan extended from smooth algebras}
\end{gathered}\\[0.4em]
\Downarrow\\[0.4em]
\text{Theorem 6.17: }(i^*\HZ^{\BL}_{p,B})^\wedge_p\simeq \HZ^{\BL}_{p,k}\\[0.4em]
\Downarrow\\[0.4em]
\text{\S6.3: field results bootstrap to mixed-characteristic Dedekind domains.}
\end{array}
\]

In particular, Theorem 6.17 is the reason Proposition 6.18 can be proved over Dedekind domains by checking after pullback to fields. This is the essential local-to-global move of the talk.


\clearpage
\section{BEM \S6.3: Voevodsky's slice conjectures over Dedekind domains and the Beilinson--Lichtenbaum equivalence}

A compressed slogan for \S6.3 is:
\[
\boxed{
\begin{gathered}
\text{field-level motivic cohomology results}\,+\,\text{base change for }\HZ^{\BL}_{p}
\\
\Longrightarrow
\\
\text{Dedekind-domain zero-slice and Beilinson--Lichtenbaum comparison.}
\end{gathered}}
\]

The section is a transfer argument: one knows the answer over fields, constructs a base-change-stable spectrum \(\HZ^{\BL}_{p,B}\), and uses it to identify the zero slice and motivic cohomology over Dedekind domains.


\subsection{Proposition 6.18: identifying \texorpdfstring{\(s^0(\1_B)^{\wedge}_p\)}{s0(1B) with HZBL}}

\begin{proposition}[BEM Proposition 6.18]
Let \(B\) be either a mixed-characteristic Dedekind domain or a field. For any prime number \(p\), the unit map
\[
       \1_B\longrightarrow \HZ^{\BL}_{p,B}
\]
induces an equivalence
\[
       s^0(\1_B)^{\wedge}_p \xrightarrow{\sim} \HZ^{\BL}_{p,B}
\]
in \(\SH(B)\).
\end{proposition}

The proposition says that the Beilinson--Lichtenbaum spectrum is the \(p\)-completed zero slice of the motivic sphere. This is the local version, over fields and Dedekind domains, of one of Voevodsky's slice expectations.


\subsubsection{Proof over fields}

Let \(B=k\) be a field. Combining Theorems 6.3 and 6.14 and passing to the inverse limit over \(m\) gives a map
\[
       \Z(\star)^A \longrightarrow \Z_p(\star)^{\BL}
\]
of \(E_\infty\)-algebras in graded presheaves of complexes on \(\Sm_k\), compatible with first Chern classes, and an equivalence modulo every power of \(p\). After applying the shearing operation of \S3.2 and then the motivic Eilenberg--MacLane functor \(H\), Remark 3.23(2) identifies the resulting map with an equivalence
\[
       s^0(\KGL_k)^{\wedge}_p \simeq \HZ^{\BL}_{p,k}.
\]
Here the shearing is the grading shift which replaces the weight-\(j\) term by its shift by \([2j]\), so that the graded cohomology theory is in the form required by \(H\).

Levine's theorem, recalled in BEM Remark 4.12, says that the unit \(\1_k\to\KGL_k\) induces an equivalence
\[
       s^0(\1_k)\xrightarrow{\sim}s^0(\KGL_k).
\]
Composing with the preceding equivalence gives
\[
       s^0(\1_k)^{\wedge}_p\xrightarrow{\sim}\HZ^{\BL}_{p,k},
\]
and this is the map induced by the unit \(\1_k\to\HZ^{\BL}_{p,k}\).

\subsubsection{Proof over mixed-characteristic Dedekind domains}

Now let \(B\) be a mixed-characteristic Dedekind domain. Since \(\HZ^{\BL}_{p,B}\) is \(p\)-complete, it is enough to identify the map after reduction modulo \(p\). Write
\[
       \HF^{\BL}_{p,B}\simeq \HZ^{\BL}_{p,B}/p.
\]
The proof isolates the two defining properties of a zero slice and the effect of the unit map.

\begin{enumerate}
  \item \(\HF^{\BL}_{p,B}\) is effective. By Proposition 3.45, effectivity can be checked after pullback to fields. For every field-valued point \(i:B\to k\), Theorem 6.17 identifies the pullback, after the harmless mod-\(p\) reduction, with \(\HF^{\BL}_{p,k}\), which is effective by the field case.

  \item \(\Fil^1_{\slic}\HF^{\BL}_{p,B}=0\). Indeed, Lemma 5.44 gives vanishing of Beilinson--Lichtenbaum cohomology in negative weights, which is exactly the vanishing of the positive effective cover needed here.

  \item Let \(F_B=\fib(\1_B\to\HZ^{\BL}_{p,B})\). The reduction \(F_B/p\) is \(1\)-effective. Again Proposition 3.45 reduces this to field-valued points; there it follows from Theorem 6.17 and the field case already proved.
\end{enumerate}

The first two points say that \(\HF^{\BL}_{p,B}\) is a zero slice, while the third says that the unit induces the same zero slice as \(\1_B/p\). Hence
\[
       s^0(\1_B)^{\wedge}_p\xrightarrow{\sim}\HZ^{\BL}_{p,B}.
\]

\subsection{Corollary 6.19: base change for the zero slice of the sphere}

\begin{corollary}[BEM Corollary 6.19]
Let \(B\) and \(k\) be either fields or mixed-characteristic Dedekind domains, and let \(i:B\to k\) be a map. Then the canonical map
\[
       i^*(s^0\1_B)\longrightarrow s^0\1_k
\]
is an equivalence in \(\SH(k)\).
\end{corollary}

By Corollary 4.52 the map is a rational equivalence. It remains to check it modulo every prime \(p\). Proposition 6.18 identifies the \(p\)-completed zero slices with \(\HZ^{\BL}_{p,B}\) and \(\HZ^{\BL}_{p,k}\), while Theorem 6.17 identifies these after pullback. Reducing the resulting comparison modulo \(p\) proves that the canonical map is a mod-\(p\) equivalence for every prime \(p\). Together with the rational equivalence, this proves Corollary 6.19.

\subsection{Proposition 6.20: from the sphere to \texorpdfstring{\(\KGL\)}{KGL}}

The next result is the Dedekind-domain version of Voevodsky's slice comparison between the motivic sphere and \(\KGL\).

\begin{proposition}[BEM Proposition 6.20]
Let \(B\) be a mixed-characteristic Dedekind domain or a field. Then:
\begin{enumerate}
  \item the map \(V_B/\beta\to s^0(V_B/\beta)\) is an equivalence;
  \item the map \(V_B\to\kgl_B\) is an equivalence;
  \item the unit map \(\1_B\to\KGL_B\) induces an equivalence
  \[
        s^0(\1_B)\xrightarrow{\sim}s^0(\KGL_B).
  \]
\end{enumerate}
\end{proposition}

Here \(V_B\) is the absolute effective motivic spectrum obtained from vector bundles with framed transfers. There is a map \(T_B\to V_B\) whose composite with \(V_B\to\KGL_B\) is the Bott element. Multiplication by this class gives a map
\[
       \beta:T_B\otimes V_B\longrightarrow V_B,
\]
and \(V_B/\beta\) denotes its cofiber. Moreover \(\KGL_B\simeq V_B[\beta^{-1}]\), and
\[
       \kgl_B:=\Fil^0_{\slic}\KGL_B
\]
is the effective cover of \(\KGL_B\).

\subsubsection{Proof of part (1)}

Over a field, the three assertions are already known: part (3) is Levine's theorem, part (2) is the field computation of \(V_k\), and then part (1) follows from
\[
       \kgl_k/\beta\simeq s^0(\kgl_k).
\]
Assume therefore that \(B\) is a mixed-characteristic Dedekind domain.

For any qcqs scheme \(X\), the unit induces an equivalence
\[
       s^0(\1_X)\xrightarrow{\sim}s^0(V_X/\beta).
\]
Indeed, Proposition 4.11 gives \(s^0(\1_X)\simeq s^0(V_X)\). From the defining cofiber sequence for \(V_X/\beta\), the cofiber of \(V_X\to V_X/\beta\) is a shift of \(T_X\otimes V_X\), hence is \(1\)-effective. Therefore \(V_X\to V_X/\beta\) does not change the zero slice.

Equivalences in \(\SH(B)\) may be checked on field-valued points. Thus it suffices to pull
\[
       V_B/\beta\longrightarrow s^0(V_B/\beta)
\]
back along every map \(i:B\to k\) to a field. The left-hand side becomes \(V_k/\beta\). For the right-hand side, the preceding comparison with \(s^0(\1)\), applied over both \(B\) and \(k\), together with Corollary 6.19 gives
\[
       i^*s^0(V_B/\beta)\simeq s^0(V_k/\beta).
\]
Hence the pullback is the canonical map \(V_k/\beta\to s^0(V_k/\beta)\), which is an equivalence by the field case. This proves part (1).

\subsubsection{Proof of part (2)}

Since \(\KGL_B\simeq V_B[\beta^{-1}]\), it suffices to show that the Bott transition maps become equivalences after taking the effective cover. For every \(j\leq0\), consider
\[
       \Fil^0_{\slic}(T_B^{\otimes j}\otimes V_B)
       \xrightarrow{\beta}
       \Fil^0_{\slic}(T_B^{\otimes(j-1)}\otimes V_B).
\]
Tensoring the defining cofiber sequence
\[
       T_B\otimes V_B\xrightarrow{\beta}V_B\longrightarrow V_B/\beta
\]
by \(T_B^{\otimes(j-1)}\) gives
\[
       T_B^{\otimes j}\otimes V_B
       \xrightarrow{\beta}
       T_B^{\otimes(j-1)}\otimes V_B
       \longrightarrow
       T_B^{\otimes(j-1)}\otimes(V_B/\beta).
\]
The functor \(\Fil^0_{\slic}\) is exact: it is a right adjoint between stable \(\infty\)-categories, and finite limits and finite colimits coincide in a stable category. Therefore the cofiber of the displayed Bott map after applying \(\Fil^0_{\slic}\) is
\[
       \Fil^0_{\slic}\bigl(T_B^{\otimes(j-1)}\otimes(V_B/\beta)\bigr).
\]
Part (1) says that \(E:=V_B/\beta\) is a zero slice. Hence \(\Fil^r_{\slic}E=0\) for every \(r\geq1\). Iterating the twist compatibility
\[
       \Fil^a_{\slic}(T_B\otimes E)
       \simeq T_B\otimes\Fil^{a-1}_{\slic}(E)
\]
from (4.7) gives
\[
       \Fil^0_{\slic}(T_B^{\otimes(j-1)}\otimes E)
       \simeq
       T_B^{\otimes(j-1)}\otimes\Fil^{1-j}_{\slic}(E).
\]
Since \(j\leq0\), one has \(1-j\geq1\), so the right-hand side vanishes. Thus every Bott transition map above is an equivalence. Passing to the colimit which inverts \(\beta\) yields
\[
       V_B\xrightarrow{\sim}\kgl_B.
\]

\subsubsection{Proof of part (3)}

Proposition 4.11 gives
\[
       s^0(\1_B)\xrightarrow{\sim}s^0(V_B).
\]
By part (2), \(V_B\simeq\kgl_B\), and the effective cover has the same zero slice as \(\KGL_B\). Therefore
\[
       s^0(\1_B)\simeq s^0(V_B)\simeq s^0(\kgl_B)\simeq s^0(\KGL_B).
\]

\begin{remark}
Proposition 6.20 is the point where the section proves the Dedekind-domain case of Voevodsky's expected comparison
\[
       s^0(\1_B)\simeq s^0(\KGL_B).
\]
The proof is indirect. It does not compute \(\KGL_B\) directly. It inserts the framed-correspondence object \(V_B\), for which base change and effectivity are better controlled.
\end{remark}

\subsection{Later consequence: Corollary 9.7}

\begin{corollary}[BEM Corollary 9.7]
Let \(X\) be any qcqs scheme, and let
\[
        f:X\longrightarrow \Spec(\Z)
\]
be the structure map. Using BEM Theorem 9.3, proved later in the paper, the unit map
\[
        \1_X\longrightarrow \KGL_X
\]
induces a commutative square of equivalences of \(E_\infty\)-algebras
\[
\begin{array}{ccc}
        f^*s^0(\1_{\Z}) & \xrightarrow{\sim} & f^*s^0(\KGL_{\Z}) \\
        \downarrow\scriptstyle{\sim} && \downarrow\scriptstyle{\sim} \\
        s^0(\1_X) & \xrightarrow{\sim} & s^0(\KGL_X).
\end{array}
\]
In particular, over every qcqs scheme \(X\), the zero slice of the motivic sphere agrees with the zero slice of \(\KGL\):
\[
        s^0(\1_X)\simeq s^0(\KGL_X).
\]
\end{corollary}

\subsubsection{What this adds beyond Proposition 6.20}

Proposition 6.20 proves
\[
        s^0(\1_B)\simeq s^0(\KGL_B)
\]
only for \(B\) a field or a mixed-characteristic Dedekind domain. Corollary 9.7, using the later Theorem 9.3, says that the same comparison holds over every qcqs scheme. Thus the section \S6.3 proves the base case over \(\Spec\Z\) and regular one-dimensional bases; the later main theorem machinery propagates this comparison to arbitrary schemes.

\subsubsection{Proof sketch}

First apply Proposition 6.20 to \(\Spec(\Z)\). This gives the top horizontal equivalence
\[
        s^0(\1_{\Z})\xrightarrow{\sim}s^0(\KGL_{\Z}).
\]
After pulling back along \(f:X\to\Spec(\Z)\), we obtain
\[
        f^*s^0(\1_{\Z})\xrightarrow{\sim}f^*s^0(\KGL_{\Z}).
\]

Next use Theorem 9.3. In the form needed here, it says that the zero slice of \(\KGL\) is stable under arbitrary base change:
\[
        f^*s^0(\KGL_{\Z})\xrightarrow{\sim}s^0(\KGL_X).
\]
This is the right vertical equivalence in the square.

It remains to explain why the bottom horizontal map
\[
        s^0(\1_X)\longrightarrow s^0(\KGL_X)
\]
is also an equivalence. Since \(\kgl_X=\Fil^0_{\slic}\KGL_X\), it is enough to show that the fibre
\[
        F_X:=\fib(\1_X\to\kgl_X)
\]
is \(1\)-effective. The reason is that if the fibre is \(1\)-effective, then it has no contribution to the zero slice, so the map \(\1_X\to\kgl_X\), and hence \(\1_X\to\KGL_X\), induces an equivalence on \(s^0\).

Now use base change. The object \(\1\) is stable under base change, and Corollary 9.6, which follows from Theorem 9.3, says that \(\kgl\) is also stable under base change:
\[
        f^*\kgl_{\Z}\simeq \kgl_X.
\]
Therefore
\[
        F_X\simeq f^*F_{\Z}.
\]
But Proposition 6.20 over \(\Spec(\Z)\) implies that \(F_{\Z}\) is \(1\)-effective. Pullback preserves effectivity, so \(F_X\) is \(1\)-effective. This proves the bottom horizontal equivalence, and hence the whole square.



\subsection{Theorem 6.21: the main theorem of \S6.3}

After the zero-slice comparison, the section turns to finite and \(p\)-adic motivic cohomology.

\begin{theorem}[BEM Theorem 6.21]
Let \(B\) be a mixed-characteristic Dedekind domain or a field, and let \(p\) be a prime. Then:
\begin{enumerate}
  \item there is a unique map of \(E_\infty\)-algebras
  \[
       s^0(\KGL_B)\longrightarrow \HZ^{\BL}_{p,B}
  \]
  in \(\SH(B)\), and after \(p\)-adic completion it is an equivalence
  \[
       s^0(\KGL_B)^{\wedge}_p\xrightarrow{\sim}\HZ^{\BL}_{p,B};
  \]

  \item every smooth \(B\)-scheme belongs to the class \(\BLp\);

  \item on \(\Sm_B\), for each \(m\geq1\), there is a unique map of \(E_\infty\)-algebras in graded presheaves of complexes
  \[
       \Z(\star)^A/p^m\longrightarrow \Z_p(\star)^{\BL}/p^m
  \]
  compatible with first Chern classes in the sense of (6.15), and this map is an equivalence.
\end{enumerate}
\end{theorem}

Thus, on smooth schemes over a field or a mixed-characteristic Dedekind domain, finite \(p\)-power \(\Aone\)-motivic cohomology is identified with Beilinson--Lichtenbaum cohomology:
\[
       \Z(j)^A/p^m\simeq \Z_p(j)^{\BL}/p^m.
\]

\subsubsection{Proof of part (1)}

Proposition 6.20(3) identifies \(s^0(\KGL_B)\) with \(s^0(\1_B)\), and Proposition 6.18 identifies the \(p\)-completion of the latter with \(\HZ^{\BL}_{p,B}\). The map is the one induced from the unit, transported through the zero-slice equivalence. This gives the unique map asserted by BEM and the equivalence
\[
       s^0(\KGL_B)^{\wedge}_p\xrightarrow{\sim}\HZ^{\BL}_{p,B}.
\]

\subsubsection{Proof of part (2): checking \texorpdfstring{\(\BLp\)}{BLp} for smooth Dedekind schemes}

The field case is Theorem 6.3, so assume that \(B\) is a mixed-characteristic Dedekind domain. One must verify conditions (A)--(H) defining \(\BLp\) for every smooth \(B\)-scheme.

Condition (H), completeness of the motivic filtration, is Remark 4.44(1). Condition (F) follows from Theorem 5.2(8). Applying \(\omega^{\infty,\mathrm{gr}}\) to part (1), and then reducing modulo \(p^m\), gives an equivalence
\[
       \Z(\star)^A/p^m \xrightarrow{\sim} \Z_p(\star)^{\BL}/p^m. \tag{6.22}
\]
of graded presheaves of complexes. At this stage one must not call (6.22) an equivalence of \(E_\infty\)-algebras: \(\omega^{\infty,\mathrm{gr}}\) is not lax symmetric monoidal in general.

The underlying equivalence of complexes transfers conditions (A)--(D) and (G) to Beilinson--Lichtenbaum cohomology. More precisely:
\begin{itemize}
  \item (A) follows from the defining truncation of \(\Z_p(j)^{\BL}\);
  \item (B) and (G) follow from BEM Remark 5.45;
  \item (C) follows from Lemma 5.44(4);
  \item under (6.22), condition (D) is reduced to the syntomic comparison condition already supplied by (F), hence to Theorem 5.2(8).
\end{itemize}

It remains to prove (E), which is the multiplicative Milnor-symbol condition. Remark 3.23(1) upgrades (6.22) to an equivalence of \(E_1\)-algebras in graded presheaves of complexes and shows that it carries
\[
       c^A_1(\calO(1))\in H_A^2(\Pone_B,\Z(1)/p^m)
\]
to
\[
       c^{\BL}_1(\calO(1))\in H_{\BL}^2(\Pone_B,\Z_p(1)/p^m).
\]
Apply Lemma 3.27 to the two first Chern class maps on \(\Z(\star)^A/p^m[2\star]\): the original map \(c_1^A\), and the map obtained by transporting \(c_1^{\BL}\) through (6.22). The hypotheses of Lemma 3.27 are supplied by Remark 4.20, the compatibility on \(\calO(1)\) just established, and Theorem 4.25. Hence the diagram
\[
  \begin{array}{ccc}
       \Gm & \xrightarrow{c^A_1} & H_A^1(-,\Z/p^m(1)) \\
       \| & & \downarrow\scriptstyle{\sim} \\
       \Gm & \xrightarrow{c^{\BL}_1} & H_{\BL}^1(-,\Z_p/p^m(1))
  \end{array}
\]
commutes.

Because (6.22) is an \(E_1\)-equivalence, it also identifies the induced bigraded ring structures. Therefore, for every smooth \(B\)-scheme \(X\), point \(x\in X\), and \(j\geq1\), the symbol map
\[
       (\calO^h_{X,x})^{\times\otimes j}
       \longrightarrow
       H_A^j(\calO^h_{X,x},\Z/p^m(j))
\]
agrees, under (6.22), with the Beilinson--Lichtenbaum symbol map. In degrees at most the weight, Beilinson--Lichtenbaum and syntomic cohomology agree; Theorem 5.42 supplies the Milnor-range calculation. This proves condition (E), hence part (2).

\subsubsection{Proof of part (3): coherent multiplicative comparison}

By part (2), the full subcategory \(\Sm_B\) lies in \(\BLp\) and is closed under \'etale maps. Theorem 6.14 therefore gives an \(E_\infty\)-map
\[
       \Z(\star)^A/p^m\longrightarrow \Z_p(\star)^{\BL}/p^m
\]
compatible with first Chern classes and proves that it is an equivalence. The uniqueness is the uniqueness built into the proof of Theorem 6.14 for any full subcategory of \(\BLp\) closed under \'etale maps.

\subsection{Corollary 6.24: Hilbert 90 and the Beilinson--Lichtenbaum equivalence}

\begin{corollary}[BEM Corollary 6.24]
Let \(B\) be a mixed-characteristic Dedekind domain or a field, and let \(j\geq0\). On \(\Sm_B\), the presheaf \(\Z(j)^A\) is Nisnevich locally supported in cohomological degrees \(\leq j\), and the canonical map
\[
       \Z(j)^A\longrightarrow L_{\Nis}\tau^{\leq j+1}L_{\et}\Z(j)^A
\]
is an equivalence.
\end{corollary}

The first assertion is an integral amplitude statement. The second is the integral Beilinson--Lichtenbaum comparison; the truncation at \(j+1\), rather than merely at \(j\), is where the Hilbert--90 boundary group enters.

\subsubsection{Proof: Nisnevich-local support}

Let \(X\) be smooth over \(B\) and \(x\in X\). By Theorem 6.21(2), \(X\in\BLp\) for every prime \(p\). Condition (A) implies that
\[
       \Z(j)^A(\calO^h_{X,x})/p
\]
is supported in cohomological degrees \(\leq j\). Condition (E) says that the degree-\(j\) mod-\(p\) classes are generated by symbols, and these symbols lift through the integral first Chern class. The long exact sequence associated with multiplication by \(p\) then shows that
\[
       H_A^n(\calO^h_{X,x},\Z(j))
\]
has no \(p\)-torsion for \(n>j\). Since this holds for every prime \(p\), these groups are torsion-free.

They are also rationally zero: Theorem 4.47 identifies rational \(\Aone\)-motivic cohomology with the appropriate Adams summands of rational homotopy \(K\)-theory, and Soulé's vanishing bound kills those summands in degrees \(n>j\) for the local rings in question. A torsion-free group which is rationally zero is zero. Thus \(\Z(j)^A\) is Nisnevich locally supported in degrees \(\leq j\).

\subsubsection{Proof: \'etale comparison and Hilbert 90}

Consider the canonical map
\[
       \Z(j)^A\longrightarrow L_{\Nis}\tau^{\leq j+1}L_{\et}\Z(j)^A. \tag{6.24}
\]
Rationally, this is an equivalence. Indeed, \(\Q(j)^A\) is locally supported in degrees \(\leq j\) by the preceding paragraph and is a direct summand of the \'etale sheaf \(\KH_{\Q}\) by Theorem 4.47(2). There is a minor sheafification subtlety: rationalization need not commute with \'etale sheafification in general. In the present situation it does, because the negative truncation of \(\Z(j)^A\) already takes rational values on \(\Sm_B\); this follows from Theorem 6.21(3) together with the vanishing of negative-degree syntomic cohomology. This is the same point isolated by BEM in their proof.

It remains to prove (6.24) modulo every prime \(p\). The key local vanishing is the Hilbert--90 statement
\[
       H^{j+1}\!\left(L_{\et}\Z(j)^A(\calO^h_{X,x})\right)=0. \tag{H90}
\]
Assuming (H90), it is enough to prove
\[
       \Z(j)^A/p
       \longrightarrow
       L_{\Nis}\tau^{\leq j}L_{\et}(\Z(j)^A/p) \tag{6.25}
\]
is an equivalence. By Theorem 6.21(3), \(\Z(j)^A/p\) may be replaced by \(\F_p(j)^{\BL}\), and (6.25) is precisely Lemma 5.44(4).

We now prove (H90). The group in (H90) is rationally zero by the rational comparison above. To prove that it is torsion-free, use the commutative square
\[
\begin{array}{ccc}
H_A^j(\calO^h_{X,x},\Z(j))
& \longrightarrow &
H^j(L_{\et}\Z(j)^A(\calO^h_{X,x})) \\
\downarrow && \downarrow \\
H_A^j(\calO^h_{X,x},\Z/p(j))
& \longrightarrow &
H^j(L_{\et}\Z(j)^A(\calO^h_{X,x})/p).
\end{array}
\]
The bottom horizontal map is an isomorphism by (6.25). The left vertical map is surjective because the degree-\(j\) classes are generated by symbols which lift integrally. Hence the right vertical map is surjective. In the long exact sequence for multiplication by \(p\) on \(L_{\et}\Z(j)^A(\calO^h_{X,x})\), this surjectivity is equivalent to the vanishing of the \(p\)-torsion in degree \(j+1\). Since this holds for every prime \(p\), the group in (H90) is torsion-free. Being also rationally zero, it vanishes. This proves (H90), hence the mod-\(p\) comparison for every \(p\), and therefore Corollary 6.24.

\section*{Reference}
T. Bachmann, E. Elmanto, M. Morrow, \emph{\(\mathbb A^1\)-invariant motivic cohomology of schemes}.\\
arXiv:2508.09915.

The numbering used in these notes follows \S\S6.2--6.3 of that paper; Corollary 9.7 is discussed only as a later consequence.

\end{document}
